\documentclass[11pt,a4paper]{article}
\usepackage[margin=25mm]{geometry}
\usepackage{amsmath,amssymb,enumitem,booktabs,array}
\usepackage[hidelinks]{hyperref}
\setlist[itemize]{nosep,leftmargin=*}
\setlist[enumerate]{leftmargin=*}
\renewcommand{\arraystretch}{1.18}
\setlength{\emergencystretch}{2em}
\newcommand{\R}{\mathbb{R}}
\newcommand{\source}[1]{\par\smallskip\noindent\textit{Source: #1.}\par\smallskip}
\title{EE6407 Genetic Algorithms and Machine Learning\\
\large Quiz 1 Preparation Notes: Evolutionary Computing Part 1 (Weeks 1--4)}
\author{Living course note}
\date{24 August 2026}
\begin{document}
\maketitle
\noindent\textbf{Scope and sources.} This is the detailed, self-contained
teaching note for Quiz~1. It covers the complete four-week Part~1 scope:
problem classification, computational complexity, evolutionary-computing
foundations, the evolutionary-algorithm loop and components, 8-queens, the
simple genetic algorithm, and optimisation context. The authoritative source
is \path{../resources/lecture-01-evolutionary-computing-part-1.pdf}, PDF
pp.~1--36, supported by the Week~1 and Week~2 transcripts. Historical
questions are used only as pattern evidence; every worked answer here is
independently derived.
\tableofcontents

\section{\enspace Evolutionary computing in one idea}
Evolutionary computing (EC) is a family of methods for searching for good solutions by maintaining and improving a population of candidates. Genetic algorithms (GAs) are an important member of this family. In lecture, the names \emph{evolutionary computing}, \emph{evolutionary algorithm}, and \emph{genetic algorithm} may be used informally, but strictly a GA is one specific type of evolutionary algorithm.

EC is useful when a problem has a large, irregular, discrete, mixed, or non-differentiable search space, or when exact optimisation would be too expensive. It does not automatically guarantee the best possible answer. It is a stochastic method for finding high-quality candidates within a computational budget.

\section{\enspace Classifying the problem before choosing a method}
Represent a system abstractly as $y=M(x)$:
\begin{itemize}
\item If $M$ and desired output $y$ are known but $x$ is unknown, we have an \textbf{optimisation} problem.
\item If input--output pairs are known but $M$ is unknown, we have a \textbf{modelling} or learning problem.
\item If $x$ and $M$ are known but $y$ is unknown, we have a \textbf{simulation} problem.
\end{itemize}
Many machine-learning tasks can be written as optimisation: select model parameters to minimise prediction error. The distinction is still useful because it clarifies what information is available and what the method must find.

The \emph{search space} $S$ contains all candidates the algorithm may consider. The \emph{feasible set} $\mathcal F\subseteq S$ contains those satisfying every hard constraint. An objective function $f:S\to\R$ measures quality. A constraint-satisfaction problem asks for any member of $\mathcal F$; a constrained optimisation problem asks for the best member.

\section{\enspace Why search can become difficult}
The number of candidate solutions can grow much faster than intuition suggests. For a route through $n$ labelled cities, there are $n!$ possible orders before considering symmetries. Exhaustive enumeration is therefore impractical even at modest $n$.

Complexity words provide context. A decision problem is in $\mathbf P$ if it can be solved in polynomial time. It is in $\mathbf{NP}$ if a proposed yes-answer can be verified in polynomial time. NP-complete problems are the hardest known problems in NP under polynomial reductions, and NP-hard problems are at least as hard but need not be decision problems. Whether $\mathbf P=\mathbf{NP}$ is unresolved. These labels motivate heuristic search; they do not prove that an EA will solve a particular instance well.

\section{\enspace The simplified evolutionary metaphor}
Evolutionary algorithms borrow four ideas from a simplified account of natural evolution: resources are limited, individuals differ, some differences affect reproductive success and are heritable, and random variation supplies new differences. In computing, \emph{fitness} is a deliberately designed numerical score rather than a biological fact.

The terms below connect the problem to its representation:
\begin{itemize}
\item A \emph{phenotype} is the candidate solution in the original problem domain.
\item A \emph{genotype} is the encoding manipulated by the algorithm.
\item A \emph{chromosome} is a complete genotype; a \emph{gene} is one component; an \emph{allele} is a possible value of that component.
\end{itemize}
The representation is an engineering decision. It should make valid solutions reachable and should work naturally with the variation operators introduced later. A poor representation can make a simple problem unnecessarily difficult.

\section{\enspace Foundations checkpoint}
\begin{itemize}
\item Can I distinguish optimisation, modelling, and simulation using $y=M(x)$?
\item Can I identify a search space, a feasible set, and an objective?
\item Can I explain why combinatorial spaces make exhaustive search difficult?
\item Can I distinguish evolutionary computing from a genetic algorithm?
\item Can I explain genotype, phenotype, chromosome, gene, and allele?
\end{itemize}
\source{\path{../resources/lecture-01-evolutionary-computing-part-1.pdf},
PDF pp.~2--16; \path{../resources/week-01-transcript.txt}.}

\clearpage
\section{Quiz 1 scope and evidence}

\subsection{What is confirmed, and what is only supporting evidence?}
This note prepares the complete four-week Part~1 scope presently contained in
\path{../resources/lecture-01-evolutionary-computing-part-1.pdf}. The supplied
deck and transcript establish a closed-book written Quiz~1 worth 10\%, held
in the final hour of Week~4. The exact date, time, and venue must still be
checked in NTULearn; tentative transcript details are not treated as
confirmed.

{\small
\begin{tabular}{@{}>{\raggedright\arraybackslash}p{0.18\textwidth}
  >{\raggedright\arraybackslash}p{0.22\textwidth}
  >{\raggedright\arraybackslash}p{0.50\textwidth}@{}}
\toprule
\textbf{Confidence} & \textbf{Evidence} & \textbf{How it is used here}\\
\midrule
Current and authoritative & Lecture 1 Part~1 deck, PDF pages 1--36 & Defines the examinable Part~1 concepts: classification, complexity, EC foundations, the EA loop and components, 8-queens, the simple GA example, EA behaviour, and optimisation context.\\
Current supporting context & Week~1 and Week~2 transcripts & Supports the assessment format and lecturer emphasis. Any tentative logistics remain unconfirmed.\\
Historical pattern only & Legacy quiz PDF pages 1--2; 2021, 2022, and 2024 walkthroughs & Suggests useful question forms involving classification, representation, encoding, fitness, and a hand-worked SGA generation. Answers below are newly derived.\\
Frequency evidence only & Exam-topic workbook, Sheet1 A1:G24 & Shows recurring older-exam themes. It does not establish the present Quiz~1 scope.\\
\bottomrule
\end{tabular}
}

\paragraph{Deliberately excluded.}
Differential evolution and particle swarm optimisation are outside this
preparation pack. Also excluded are advanced crossover families, schema
theory, niching, multi-objective optimisation, and machine learning. These
topics appeared in older or later materials but are not in the approved
Part~1 scope.

\source{Lecture~1 Part~1 deck, PDF pp.~1--36; Week~1 and Week~2 transcripts;
historical-pattern sources indexed in \path{../STATUS.md}.}

\clearpage
\section{One-page memory core}
\begin{tabular}{@{}>{\bfseries\raggedright\arraybackslash}p{0.23\textwidth}
  >{\raggedright\arraybackslash}p{0.71\textwidth}@{}}
\toprule
Problem classification & For $y=M(x)$: unknown $x$ means optimisation; unknown $M$ means modelling/learning; unknown $y$ means simulation.\\
Search language & Search space $S$: all encodable candidates. Feasible set $\mathcal F\subseteq S$: candidates satisfying hard constraints. Objective $f$: quality. CSP seeks any feasible candidate; COP seeks the best feasible candidate.\\
Complexity & $\mathbf P$: solvable in polynomial time. $\mathbf{NP}$: a proposed yes-certificate is verifiable in polynomial time. NP-complete: in NP and NP-hard. NP-hard: at least as hard as every NP problem, but may be an optimisation problem and need not be in NP.\\
EA cycle & Initialise $\rightarrow$ evaluate $\rightarrow$ select parents $\rightarrow$ vary by recombination/mutation $\rightarrow$ evaluate offspring $\rightarrow$ select survivors $\rightarrow$ test termination.\\
Representation & Phenotype is the domain solution; genotype is its encoding. A useful encoding makes intended solutions reachable and lets operators preserve or repair validity.\\
Selection & Fitness-proportionate selection: $p_i=f_i/\sum_j f_j$ when fitnesses are non-negative and larger is better. Selection favours; it does not guarantee.\\
8-queens & Permutation $q=(q_1,\ldots,q_N)$ places one queen in column $i$, row $q_i$. Queens $i,j$ conflict diagonally when $|q_i-q_j|=|i-j|$. With $C$ diagonal conflicts, $f=\binom N2-C$.\\
Binary decoding & For bits $b_{L-1}\cdots b_0$, $z=\sum_{k=0}^{L-1}2^k b_k$. Optional interval map: $x=a+\dfrac{z}{2^L-1}(b-a)$.\\
Mutation checks & Across $N$ strings of length $L$, expected flipped bits are $NLp_m$. Probability of no flip is $(1-p_m)^{NL}$.\\
Optimisation context & A neighbourhood defines local moves. A local optimum is best only in its neighbourhood; a global optimum is best in the whole feasible set. No single optimiser is best across all possible problems.\\
\bottomrule
\end{tabular}

\paragraph{Closed-book recall test.}
Cover the table and reconstruct the three $y=M(x)$ cases, the P/NP relationships, the seven EA stages, the 8-queens conflict test, roulette selection, and binary decoding.

\clearpage
\section{Five worked archetypes}

\subsection{Archetype 1: classify the problem before choosing an algorithm}
\textbf{Question.} A factory knows a simulator $M$ that predicts energy use $y$ for a machine setting $x$. Classify each task:
\begin{enumerate}[label=(\alph*)]
\item Find a setting that minimises energy while meeting production constraints.
\item Predict the energy used by one specified setting.
\item Infer a predictive relationship from logged settings and measured energy.
\end{enumerate}
\textbf{Construction.} Write $y=M(x)$ and ask which object is unknown.
\begin{enumerate}[label=(\alph*)]
\item $M$ is known and the desired low-energy outcome is specified, but $x$ is unknown: \textbf{optimisation}. The search space contains encodable settings, the feasible set satisfies production constraints, and the objective is predicted energy.
\item $x$ and $M$ are known, while $y$ is requested: \textbf{simulation}.
\item Input--output examples are known, while $M$ is unknown: \textbf{modelling/learning}.
\end{enumerate}
\textbf{Trap.} The use of a computer model does not automatically make a task simulation. Classification depends on what is unknown.

\subsection{Archetype 2: reason correctly about P, NP, NP-complete, and NP-hard}
\textbf{Question.} Assess the claim that NP means non-polynomial and therefore an NP-hard optimisation problem cannot be checked efficiently.

\textbf{Answer.} The claim combines two errors.
\begin{enumerate}
\item NP means \emph{nondeterministic polynomial time}, commonly characterised by polynomial-time verification of a proposed yes-certificate for a decision problem. It does not mean that the problem is proved to require non-polynomial time.
\item NP-hard is a hardness relation. An NP-hard problem need not itself be a decision problem or belong to NP. For an optimisation problem, checking that a candidate is feasible and calculating its objective may be easy even when proving that it is globally optimal is hard.
\end{enumerate}
If a decision problem is both in NP and NP-hard, it is NP-complete. Every problem in P is also in NP. Whether $\mathbf P=\mathbf{NP}$ remains unresolved.

\subsection{Archetype 3: trace the evolutionary algorithm loop}
\textbf{Question.} A population contains four binary strings. Describe one generation without confusing parent selection and survivor selection.
\begin{enumerate}
\item \textbf{Represent and initialise:} choose an encoding and create the initial population.
\item \textbf{Evaluate:} decode each genotype, calculate its objective, and convert it to the fitness used by selection.
\item \textbf{Select parents:} sample candidates for reproduction. Better fitness normally means greater probability, not certainty.
\item \textbf{Recombine:} exchange genetic material between selected parents to make offspring.
\item \textbf{Mutate:} randomly change genes, supplying new variation.
\item \textbf{Evaluate offspring:} decode and score the new candidates.
\item \textbf{Select survivors:} decide which parents and/or offspring form the next population.
\item \textbf{Terminate or repeat:} stop at a target, budget, or stagnation condition; otherwise begin the next generation.
\end{enumerate}
\textbf{Trap.} Parent selection answers who reproduces; survivor selection answers who remains. They can use different rules.

\subsection{Archetype 4: represent and score 8-queens}
For $N=4$, let $q_i$ be the row of the queen in column $i$. The permutation $q=(2,4,1,3)$ automatically gives one queen in each row and column. For every pair $i<j$, compare $|q_i-q_j|$ with $|i-j|$:
\[
\begin{array}{c|cccccc}
(i,j)&(1,2)&(1,3)&(1,4)&(2,3)&(2,4)&(3,4)\\
\hline
|q_i-q_j|&2&1&1&3&1&2\\
|i-j|&1&2&3&1&2&1
\end{array}
\]
No pair is equal, so $C=0$ and $f(q)=\binom42-C=6$. A swap mutation exchanging the last two positions gives $q'=(2,4,3,1)$. It remains a valid permutation, but pair $(2,3)$ conflicts, so $C=1$ and $f(q')=5$.

\textbf{Operator lesson.} Swap mutation preserves the permutation. Recombination must likewise preserve it or be followed by repair.

\subsection{Archetype 5: calculate one complete SGA generation}
Maximise $f(x)=x^2$ for an unsigned four-bit integer. Start with:
\[
\begin{array}{c|c|c|c}
\text{chromosome}&x&f(x)&p_i=f_i/\sum f_j\\
\hline
0011&3&9&9/311=0.0289\\
0101&5&25&25/311=0.0804\\
1001&9&81&81/311=0.2605\\
1110&14&196&196/311=0.6302
\end{array}
\]
The fitness sum is $311$. Suppose roulette draws produce the mating pool
$[1110,1001,1110,0101]$. Pair adjacent parents and cross after bit 2:
\[
\begin{aligned}
11|10,\;10|01&\longrightarrow 11|01=1101,\quad 10|10=1010,\\
11|10,\;01|01&\longrightarrow 11|01=1101,\quad 01|10=0110.
\end{aligned}
\]
Suppose mutation flips the final bit of the fourth offspring only, so $0110\to0111$. With generational replacement the next population is
$[1101,1010,1101,0111]$, decoding to $[13,10,13,7]$, with fitnesses $[169,100,169,49]$. Its total fitness is $487$ and its best fitness is $169$.

\textbf{Checks.} The roulette outcomes and mutation event are sampled events, not guarantees. Decode after variation, recompute fitness, and use the correct generation's fitness sum.

\clearpage
\section{Progressive drills}
\textit{Attempt all twelve without notes. Hints are brief; complete solutions begin on the next page.}
\begin{enumerate}[label=\textbf{D\arabic*.}]
\item For $y=M(x)$, classify: (i) estimating $M$ from examples; (ii) finding $x$ for a desired $y$ with known $M$; (iii) finding $y$ for known $x,M$.\par\emph{Hint: identify the unknown object.}
\item A delivery problem has all possible route encodings $S$, a maximum-distance rule, and total fuel use $f$. Identify the search space, feasible set, and objective.\par\emph{Hint: constraints define a subset of the search space.}
\item How many ordered routes are there through eight labelled cities before symmetry reductions? Why does this number matter?\par\emph{Hint: order all eight objects.}
\item Correct the claim that every NP-hard problem is an NP-complete decision problem and NP means non-polynomial.\par\emph{Hint: separate membership from hardness.}
\item Put these in a valid EA order: mutation, termination test, initialisation, offspring evaluation, parent selection, survivor selection, population evaluation, recombination.\par\emph{Hint: candidates must exist before evaluation.}
\item Three candidates have fitnesses $2,3,5$. Calculate their roulette-selection probabilities.\par\emph{Hint: divide by total fitness.}
\item Explain in one sentence why repeatedly selecting only the current best candidate can be risky.\par\emph{Hint: think about diversity and local optima.}
\item For $q=(2,4,1,3)$ in 4-queens, calculate the diagonal conflicts and fitness $\binom42-C$.\par\emph{Hint: test every pair.}
\item Parents $(1,2,3,4)$ and $(4,3,2,1)$ undergo ordinary one-point crossover after position 2. Write the children and explain the validity problem.\par\emph{Hint: a permutation contains each row once.}
\item Decode the unsigned binary string $10110$.\par\emph{Hint: use place values $16,8,4,2,1$.}
\item Map $10110$ linearly to $[10,20]$, including both endpoints.\par\emph{Hint: use $10+\frac{z}{31}(20-10)$.}
\item Six strings each have length eight and independent bit-mutation probability $p_m=0.005$. Find the expected flips and probability of no flips.\par\emph{Hint: there are $48$ trials.}
\end{enumerate}

\clearpage
\section{Solutions to the progressive drills}
\begin{enumerate}[label=\textbf{D\arabic*.}]
\item (i) Modelling/learning; (ii) optimisation; (iii) simulation.
\item $S$ is every encodable route; $\mathcal F$ is the subset obeying the maximum-distance rule; minimise total fuel use $f$.
\item $8!=40320$. Factorial growth makes exhaustive enumeration rapidly expensive.
\item NP is polynomial-time verification of yes-certificates for decision problems, not non-polynomial. NP-hard problems need not be in NP or be decision problems. NP-complete problems are both in NP and NP-hard.
\item
\[
\begin{aligned}
\text{initialise}&\to\text{evaluate population}\to\text{select parents}
\to\text{recombine}\to\text{mutate}\\
&\to\text{evaluate offspring}\to\text{select survivors}
\to\text{test termination}.
\end{aligned}
\]
\item Total $10$; probabilities $0.2,0.3,0.5$, summing to $1$.
\item It can destroy diversity and cause premature convergence around a local optimum.
\item No conflicting pair, so $C=0$ and $f=6$.
\item Children $(1,2,2,1)$ and $(4,3,3,4)$ repeat and omit rows; use a permutation-preserving crossover or repair.
\item $10110_2=16+4+2=22$.
\item $x=10+\frac{22}{31}(10)=17.0968$ approximately.
\item Expected flips $=6(8)(0.005)=0.24$; no-flip probability $=0.995^{48}\approx0.7862$.
\end{enumerate}

\clearpage
\section{Mock Quiz A: foundations and a hand-worked generation}
\noindent\textbf{Conditions:} 60 minutes, closed book, 100 marks.

\subsection*{Question 1: Problem formulation (20 marks; 10 minutes)}
A trained surrogate predicts drag $y=M(x)$ from an airfoil design $x$.
\begin{enumerate}[label=(\alph*)]
\item Classify finding the lowest-drag design subject to strength constraints. Identify $S$, $\mathcal F$, and the objective. \hfill[10]
\item Classify (i) predicting drag for a given design and (ii) learning the surrogate from examples. \hfill[6]
\item Distinguish a CSP from a constrained optimisation problem. \hfill[4]
\end{enumerate}

\subsection*{Question 2: Complexity and EC foundations (20 marks; 10 minutes)}
\begin{enumerate}[label=(\alph*)]
\item Define P, NP, NP-complete, and NP-hard. \hfill[12]
\item Explain why NP-hardness motivates heuristic search but does not prove an EA will perform well. \hfill[4]
\item Distinguish genotype from phenotype, with one example. \hfill[4]
\end{enumerate}

\subsection*{Question 3: The EA loop (20 marks; 10 minutes)}
\begin{enumerate}[label=(\alph*)]
\item Write one complete generation of a general EA in a correct order. \hfill[12]
\item Distinguish parent selection from survivor selection. \hfill[4]
\item Give two sensible termination conditions. \hfill[4]
\end{enumerate}

\subsection*{Question 4: 4-queens representation (20 marks; 12 minutes)}
Use $q_i=$ row in column $i$ and fitness $f=\binom42-C$.
\begin{enumerate}[label=(\alph*)]
\item For $q=(1,3,4,2)$, list every diagonal conflict and calculate $f$. \hfill[10]
\item Explain which constraints the permutation enforces automatically. \hfill[4]
\item Explain why ordinary one-point crossover may fail and give a remedy. \hfill[6]
\end{enumerate}

\subsection*{Question 5: Simple GA calculation (20 marks; 18 minutes)}
Maximise $f(x)=x^2$ using three-bit unsigned strings. The population is $010,011,101,110$.
\begin{enumerate}[label=(\alph*)]
\item Decode, calculate all fitnesses, and calculate roulette probabilities. \hfill[10]
\item The mating pool is $110,101,011,110$. Cross adjacent parents after bit 1. \hfill[6]
\item Flip the last bit of the first child only. Give the next population and best fitness. \hfill[4]
\end{enumerate}

\clearpage
\section{Mock Quiz A: marking guide}
\begin{description}[leftmargin=0pt,labelindent=0pt]
\item[Q1 (20).] Optimisation (2); $S$ all encodable designs (2); $\mathcal F$ strength-feasible designs (2); minimise drag (2); known $M$/unknown $x$ (2). Simulation (3), modelling/learning (3). CSP seeks any feasible candidate (2); constrained optimisation seeks the best (2).
\item[Q2 (20).] P solvable in polynomial time (3); NP has polynomially verifiable yes-certificates (3); NP-complete is in NP and NP-hard (3); NP-hard need not be in NP (3). Motivation without algorithm-specific guarantee (4). Genotype is encoding; phenotype is domain solution, with example (4).
\item[Q3 (20).] Initialise, evaluate, parent select, recombine, mutate, evaluate offspring, survivor select, terminate/repeat: 12. Parent selection chooses reproducers; survivor selection forms the next population: 4. Any two of target, evaluation/time budget, or stagnation: 4.
\item[Q4 (20).] Only $(2,3)$ conflicts (6); $C=1$, $f=5$ (4). One queen per column and row (4). Ordinary crossover can duplicate/omit rows (3); permutation-preserving crossover or repair (3).
\item[Q5 (20).] Decodings $2$, $3$, $5$, $6$ (2); fitnesses $4$, $9$,
$25$, $36$ (2); sum $74$ (2); probabilities $4/74$, $9/74$, $25/74$,
$36/74$ (4). The crosses are
\[
1|10,\;1|01\to101,\;110
\qquad\text{and}\qquad
0|11,\;1|10\to010,\;111
\]
(6). Mutate $101\to100$; the next population is $100$, $110$, $010$,
$111$ (2), and its best fitness is $49$ (2).
\end{description}
\noindent\textbf{Total: 100 marks.}

\clearpage
\section{Mock Quiz B: representation, operators, and optimisation context}
\noindent\textbf{Conditions:} 60 minutes, closed book, 100 marks.

\subsection*{Question 1: Representation and encoding (20 marks; 12 minutes)}
\begin{enumerate}[label=(\alph*)]
\item Define chromosome, gene, allele, genotype, and phenotype. \hfill[10]
\item Decode $11010_2$ and map it linearly to $[-2,3]$, including endpoints. \hfill[6]
\item State two properties of a useful representation. \hfill[4]
\end{enumerate}

\subsection*{Question 2: Selection and variation (20 marks; 10 minutes)}
Four candidates have fitnesses $1,2,3,4$.
\begin{enumerate}[label=(\alph*)]
\item Calculate roulette probabilities. \hfill[6]
\item Explain why the best candidate is not guaranteed to be selected. \hfill[4]
\item Distinguish recombination from mutation and state one role of each. \hfill[6]
\item State one danger of excessive selection pressure. \hfill[4]
\end{enumerate}

\subsection*{Question 3: 5-queens fitness (20 marks; 12 minutes)}
Use $q=(1,2,3,5,4)$ and fitness $f=\binom52-C$.
\begin{enumerate}[label=(\alph*)]
\item Identify all diagonal conflicts. \hfill[12]
\item Calculate $C$ and $f$. \hfill[4]
\item Explain why swap mutation preserves row/column feasibility. \hfill[4]
\end{enumerate}

\subsection*{Question 4: SGA checks (25 marks; 16 minutes)}
Five strings have length six; bit mutation is independent with $p_m=0.01$.
\begin{enumerate}[label=(\alph*)]
\item Find expected flips and probability of no flips. \hfill[8]
\item Cross $101|011$ and $010|110$ at the shown point. \hfill[5]
\item Decode both children as unsigned integers. \hfill[4]
\item For minimising $g(x)=(x-20)^2$, explain why raw $g$ is wrong in a larger-is-better roulette rule and give a valid transformation. \hfill[8]
\end{enumerate}

\subsection*{Question 5: Search behaviour (15 marks; 10 minutes)}
\begin{enumerate}[label=(\alph*)]
\item Distinguish local and global optima using a neighbourhood. \hfill[6]
\item Explain exploration versus exploitation. \hfill[5]
\item State the practical message of no free lunch. \hfill[4]
\end{enumerate}

\clearpage
\section{Mock Quiz B: marking guide}
\begin{description}[leftmargin=0pt,labelindent=0pt]
\item[Q1 (20).] Chromosome/genotype, gene/component, allele/value, genotype/encoding, phenotype/domain solution: 2 each. $11010_2=26$ (2); $x=-2+\frac{26}{31}(5)=2.1935$ (4). Any two useful representation properties: 2 each.
\item[Q2 (20).] Total $10$, probabilities $0.1,0.2,0.3,0.4$ (6). Stochastic selection is not certainty (4). Recombination combines parental material (3); mutation supplies random new variation (3). Excess pressure can destroy diversity and cause premature convergence (4).
\item[Q3 (20).] Conflicts $(1,2),(1,3),(2,3),(4,5)$: 3 each. Thus $C=4$, $f=10-4=6$ (4). Swapping retains each row once, while columns are positions (4).
\item[Q4 (25).] There are $30$ bits; expected flips $0.3$ (4); no-flip probability $0.99^{30}\approx0.7397$ (4). Children $101110$ and $010011$ (5), decoding to $46$ and $19$ (4). Raw $g$ rewards larger error (4); for example $F=1/(1+g)$ is positive and decreases with error (4).
\item[Q5 (15).] Local is best in its neighbourhood; global is best in the whole feasible set (6). Exploration samples new regions; exploitation improves promising regions (5). No optimiser dominates all possible problems, so problem knowledge and empirical testing matter (4).
\end{description}
\noindent\textbf{Total: 100 marks.}

\clearpage
\section{Eight-day preparation schedule: 24--31 August}
\begin{tabular}{@{}>{\bfseries\raggedright\arraybackslash}p{0.12\textwidth}
  >{\raggedright\arraybackslash}p{0.27\textwidth}
  >{\raggedright\arraybackslash}p{0.55\textwidth}@{}}
\toprule
Date & Focus & Minimum deliverable\\
\midrule
24 Aug & Scope and classification & Reproduce the scope map; solve D1--D3 without notes.\\
25 Aug & Complexity and EC foundations & Explain P/NP/NP-complete/NP-hard aloud; solve D4.\\
26 Aug & EA loop and components & Draw the complete loop from memory; solve D5--D7.\\
27 Aug & Representation and 8-queens & Recalculate Archetype 4; solve D8--D9.\\
28 Aug & SGA arithmetic & Recalculate Archetype 5; solve D10--D12 and check every sum.\\
29 Aug & Timed Mock A & Attempt in 60 minutes; mark in a different colour; relearn every lost mark.\\
30 Aug & Timed Mock B & Attempt in 60 minutes; mark and write a one-line cause for each error.\\
31 Aug & Retrieval and weak spots & Recite the memory core, redo missed questions, and stop heavy study early.\\
\bottomrule
\end{tabular}

\subsection{Quiz-day checklist}
\begin{itemize}
\item Verify the official date, time, venue, and permitted materials in NTULearn.
\item Before calculating, write the encoding, objective direction, and fitness convention.
\item For roulette selection, calculate the fitness sum and check probabilities sum to 1.
\item After crossover or mutation, check chromosome length and representation validity.
\item For 8-queens, test each pair once using $i<j$; do not double-count.
\item For complexity questions, do not expand NP as non-polynomial.
\item Reserve the final five minutes for arithmetic and unanswered subparts.
\end{itemize}

\source{\path{../resources/lecture-01-evolutionary-computing-part-1.pdf},
PDF pp.~1--36. Historical sources were used only to design original question
patterns; all calculations and solutions here were independently derived.}
\end{document}
